\documentclass[10pt,a4paper]{article}
\usepackage[utf8]{inputenc}
\usepackage[ngerman]{babel}
\usepackage[T1]{fontenc}
\usepackage{amsmath}
\usepackage{amsfonts}
\usepackage{amssymb}
\usepackage{graphicx}
\usepackage{lmodern}
\usepackage{physics}
\usepackage[left=1cm,right=1cm,top=2cm,bottom=1.5cm]{geometry}
\usepackage{siunitx}
\usepackage{fancyhdr}
\usepackage{enumerate}
\usepackage{mhchem}
\usepackage{mathtools}
\usepackage{graphicx}
\usepackage{float}
\usepackage{xcolor}
\usepackage{mdframed}
\usepackage{csquotes}
\usepackage{trfsigns}
\usepackage{capt-of}


\usepackage{amsmath}
\usepackage{graphicx}
\usepackage{enumitem}
\usepackage{geometry}
\geometry{margin=2.5cm}


\sisetup{locale=DE}
\sisetup{per-mode = symbol-or-fraction}
\sisetup{separate-uncertainty=true}
\DeclareSIUnit\year{a}
\DeclareSIUnit\clight{c}
\mdfdefinestyle{exercise}{
	backgroundcolor=black!10,roundcorner=8pt,hidealllines=true,nobreak
}

\begin{document}
	\twocolumn
	\pagestyle{fancy}
	\lhead{Formelsammlung Messtechnik}
	\rhead{\today}
	\section{Messabweichung}
		\begin{mdframed}[style=exercise]
			\begin{align}
					\Delta y = \sum_{i=1}^{n} \frac{\partial f}{\partial x_i} \, \Delta x_i \quad \text{für} \quad \Delta x_i \ll x_i
			\end{align}
			$\delta$y = Gesamtunsicherheit \par
			$\frac{\partial f}{\partial x_i}$ = Empfindlichkeit
		\end{mdframed}
	
	\section{Empfindlichkeit}
		\begin{mdframed}[style=exercise]
			\begin{align}
				E &= \left. \frac{dy}{dx} \right|_{\text{AP}} \approx \left. \frac{\Delta y}{\Delta x} \right|_{\text{AP}} \\
				y_m &\approx y_{\text{AP}} + E \cdot (x_m - x_{\text{AP}})
			\end{align}
		\end{mdframed}
		→ Linearisierung um den AP
		→ Zulässiger Lin.-Fehler bestimmt Größe des Bereichs um den AP
		Hinweis: Falls Gerade y = m $\times$ x → E = m = Steigung = konstant
		
	\section{Empfindlichkeit im linearen Bereich}
		\begin{mdframed}[style=exercise]
			\begin{align}
			E &= \frac{dy}{dx} \approx \frac{\Delta y}{\Delta x} \\
			y_m &\approx y_0 + \frac{\Delta y}{\Delta x} (x_m - x_0) \\
			&= y_0 + E(x_m - x_0)
			\end{align}
		\end{mdframed}
		Im oberen Beispiel wird nicht symmetrisch um den AP, sondern entlang eines Eingangsbereichs linearisiert.
	

	
	\pagebreak
	
	
	\section{Platzhalter}
	\begin{mdframed}[style=exercise]
		\begin{align}
			\Delta y = \sum_{i=1}^{n} \frac{\partial f}{\partial x_i} \, \Delta x_i \quad \text{für} \quad \Delta x_i \ll x_i
		\end{align}
		$\delta$y = Gesamtunsicherheit \par
		$\frac{\partial f}{\partial x_i}$ = Empfindlichkeit
	\end{mdframed}
	
	\pagebreak
	
	
	\section{Platzhalter}
	\begin{mdframed}[style=exercise]
		\begin{align}
			\Delta y = \sum_{i=1}^{n} \frac{\partial f}{\partial x_i} \, \Delta x_i \quad \text{für} \quad \Delta x_i \ll x_i
		\end{align}
		$\delta$y = Gesamtunsicherheit \par
		$\frac{\partial f}{\partial x_i}$ = Empfindlichkeit
	\end{mdframed}
\end{document}